% Paste this block into sections/5_experiments.tex after the current
% computation-communication overlap analysis. It uses only packages already
% loaded in paper.tex: graphicx is enough for \resizebox.

\subsection{Supplemental Validation}
\label{sec:supplemental_validation}

To address correctness, overhead, scalability, and recent-method comparison
concerns, we add a controlled supplemental evaluation. Unless otherwise stated,
the controlled system harness uses 16 MoE layers, 64 routed experts, Top-$K{=}6$,
a 512-entry expert cache, and 8\,MB expert transfers. All baselines share the
same cache capacity and host-to-device transfer model. The proxy rows for recent
systems are same-harness approximations used to isolate design choices; they are
not official reproductions of the corresponding systems.

\textbf{Correctness of verified prefetching.}
SPICE changes expert residency timing but not target expert execution. In a
verified MoE execution harness, on-demand offloading and SPICE produce identical
logits: the maximum logit difference is 0.0 and the exact argmax match rate is
100\%. The baseline pseudo-perplexity and SPICE verified pseudo-perplexity are
both 1035.21. This supports the invariant in Section~\ref{sec:methodology}:
fallback traffic affects latency, not model semantics.

\begin{table}[t]
\centering
\caption{Controlled same-harness offloading and prefetch comparison.}
\label{tab:supp_prefetch}
\scriptsize
\resizebox{\linewidth}{!}{
\begin{tabular}{lccccc}
\hline
Policy & TPOT (ms) & Hit & Fallback & H2D (GB) & PCIe active \\
\hline
Naive & 71.25 & 0.00\% & 100.00\% & 384.00 & 43.86\% \\
LRU & 44.67 & 85.04\% & 14.96\% & 57.44 & 10.46\% \\
MoE-Offloading & 44.67 & 85.04\% & 14.96\% & 57.44 & 10.46\% \\
Pre-gated & 41.83 & 94.15\% & 5.85\% & 147.27 & 28.65\% \\
SPICE & 41.03 & 99.76\% & 0.24\% & 198.80 & 39.43\% \\
\hline
\end{tabular}}
\end{table}

Table~\ref{tab:supp_prefetch} shows that SPICE reduces critical-path fallback
traffic from 5.85\% for Pre-gated to 0.24\%. Its H2D volume is higher because it
prefetches more aggressively, but most of this traffic is overlapped with target
execution rather than paid on the critical path.

\begin{table}[t]
\centering
\caption{Draft-path overhead and LoRE ablation.}
\label{tab:supp_overhead}
\scriptsize
\resizebox{\linewidth}{!}{
\begin{tabular}{lcccc}
\hline
Variant & TPOT (ms) & Fallback & Draft overhead & Online overhead \\
\hline
SPICE offline & 41.04 & 0.24\% & 491.52 ms & 0.00 ms \\
SPICE online & 42.96 & 0.24\% & 491.52 ms & 983.04 ms \\
SPICE without LoRE & 42.56 & 5.10\% & 491.52 ms & 0.00 ms \\
\hline
\end{tabular}}
\end{table}

Table~\ref{tab:supp_overhead} separates the offline draft path from online
self-correction. Online adaptation is measurable rather than free: enabling
online updates increases TPOT from 41.04\,ms to 42.96\,ms in this setup. We
therefore use the offline draft path as the default latency configuration and
treat online adaptation as an optional robustness mode. Removing LoRE increases
fallback from 0.24\% to 5.10\%, showing that the low-rank state correction is
responsible for much of the verified prefetch accuracy.

\begin{table}[t]
\centering
\caption{Top-$K$ stress study under a tighter cache budget.}
\label{tab:supp_topk}
\scriptsize
\resizebox{\linewidth}{!}{
\begin{tabular}{ccccc}
\hline
Top-$K$ & TPOT (ms) & Hit & Fallback & PCIe active \\
\hline
2 & 41.03 & 99.37\% & 0.63\% & 22.95\% \\
4 & 45.78 & 76.84\% & 23.16\% & 51.05\% \\
6 & 53.15 & 60.98\% & 39.02\% & 79.76\% \\
8 & 61.82 & 49.95\% & 50.05\% & 100.00\% \\
10 & 71.34 & 41.67\% & 58.33\% & 100.00\% \\
12 & 78.80 & 39.46\% & 60.54\% & 100.00\% \\
\hline
\end{tabular}}
\end{table}

Table~\ref{tab:supp_topk} clarifies SPICE's scalability boundary. As Top-$K$
increases, the required expert-residency working set grows and PCIe activity
reaches saturation by $K{=}8$. SPICE therefore improves the offloaded regime by
using available overlap windows more effectively, but it does not remove the
physical bandwidth limit of the host-to-device link.

\begin{table}[t]
\centering
\caption{Controlled proxy comparison with recent MoE inference systems.}
\label{tab:supp_proxy}
\scriptsize
\resizebox{\linewidth}{!}{
\begin{tabular}{lccc}
\hline
Variant & TPOT (ms) & Hit & Fallback \\
\hline
ExpertFlow proxy & 41.70 & 98.66\% & 1.34\% \\
AdapMoE proxy & 42.08 & 94.38\% & 5.62\% \\
SP-MoE proxy & 43.04 & 95.39\% & 4.61\% \\
MoE-SpeQ proxy & 41.98 & 95.72\% & 4.28\% \\
SPICE verified & 41.03 & 99.78\% & 0.22\% \\
\hline
\end{tabular}}
\end{table}

Table~\ref{tab:supp_proxy} compares SPICE against controlled proxy variants of
recent systems. The result is not intended as an official reproduction; rather,
it isolates the effect of routing-only verified prefetching under one cache and
PCIe model. SPICE has the lowest fallback rate because its speculative decision
controls expert residency and is verified against target routing before expert
execution.

\textbf{Power and timeline telemetry.}
We also collect GPU telemetry to avoid inferring energy from latency alone. In a
copy-and-compute workload, the measured GPU power averages 184.32\,W and peaks
at 452.98\,W over a 44.90\,s sampling window. The workload transfers 16\,GB in
0.569\,s, corresponding to an effective 28.13\,GB/s. We also collect an Nsight
Systems CUDA trace for the same workload. These measurements are telemetry
evidence for PCIe overlap and runtime overhead; we do not claim full
system-level energy savings without paired baseline-vs-SPICE energy runs.

